(1) 相干态的数态展开由湮灭算符作用 \(a|n\rangle=\sqrt n|n-1\rangle\) 递推得到 \(|\alpha\rangle=C\sum_{n=0}^\infty\alpha^n|n\rangle/\sqrt{n!}\)。归一化条件 \(1=|C|^2\sum_{n\ge0}|\alpha|^{2n}/n!=|C|^2e^{|\alpha|^2}\)，故取 \(C=e^{-|\alpha|^2/2}\)，得到规范化的标准相干态表达式。
(2) 测得 \(n\) 的概率为 \(P_n=|\langle n|\alpha\rangle|^2=e^{-|\alpha|^2}|\alpha|^{2n}/n!\)，即参数 \(\lambda=|\alpha|^2\) 的泊松分布。由指数母函数可得 \(\langle n\rangle=|\alpha|^2\)、\(\operatorname{Var}(n)=|\alpha|^2\)，并且 Fano 因子为一；这是相干态的数涨落而非固定数态的零方差。
(3) 在自由演化 \(e^{-iHt/\hbar}\) 下，数态第 \(n\) 项获得相位 \(e^{-i\omega(n+1/2)t}\)，整体公共相位 \(e^{-i\omega t/2}\) 可去掉，相干参数变为 \(\alpha(t)=\alpha(0)e^{-i\omega t}\)。由 \(x=\sqrt{\hbar/(2m\omega)}(a+a^\dagger)\)，得 \(\langle x(t)\rangle=\sqrt{2\hbar/(m\omega)}\,\operatorname{Re}[\alpha(0)e^{-i\omega t}]\)，正是经典谐振频率的正弦运动。
(4) 用 \(p=-i\sqrt{m\hbar\omega/2}(a-a^\dagger)\) 和相干态算符平均，有 \((\Delta x)^2=\hbar/(2m\omega)\)、\((\Delta p)^2=m\hbar\omega/2\)，从而 \(\Delta x\Delta p=\hbar/2\)。二者在任意时间恒定，意味着概率包络整体随经典轨迹平移而不展宽；这是谐振子的等间隔能级与相干态特殊相位结构共同造成，不代表任意量子波包均无扩散。
